\section{Conclusion, Limitations, and Future Research Directions}

This paper proposed a fractional stochastic transportation model integrating Caputo long-memory dynamics, stochastic demand, carbon trading, and CVaR-based risk aversion within a unified two-stage optimization framework. Unlike conventional transportation formulations that rely on memoryless balance equations, the proposed model explicitly accounts for the cumulative influence of past operational states through a Caputo fractional balance equation. Consequently, transportation decisions are affected not only by current conditions but also by the historical evolution of congestion, inventory, service performance, and operational disturbances.

The proposed formulation combines strategic first-stage transportation decisions with scenario-dependent recourse actions, allowing shipment quantities, shortages, inventory-service states, and carbon-trading decisions to adapt to realized demand. The fractional state equation is transformed into a tractable deterministic equivalent through the classical L1 discretization, while the stochastic structure is addressed using Sample Average Approximation (SAA) combined with Progressive Hedging (PH). In addition, several theoretical properties were established, including the existence of an optimal solution, the impossibility of simultaneous carbon-allowance buying and selling under a positive bid-ask spread, and the stability of the optimal solution under moderate scenario perturbations.

The computational study demonstrates that the fractional order significantly influences transportation behavior. Stronger memory effects encourage more conservative planning by anticipating shipments and maintaining higher effective inventory-service levels. Increasing demand volatility raises both expected transportation cost and tail risk, whereas tighter carbon allowances increase the reliance on carbon trading and encourage cleaner transportation decisions. Higher levels of risk aversion improve system robustness by reducing extreme shortage costs, although the corresponding economic benefit gradually diminishes. These results highlight that explicitly modeling memory can substantially modify both operational and environmental decisions compared with classical first-order transportation models.

Despite these contributions, several limitations remain. First, the computational experiments rely on synthetic benchmark instances designed to isolate the effects of the principal model parameters under controlled conditions. Although this approach provides a transparent methodological assessment, it cannot replace validation using real transportation, inventory, congestion, and carbon-emission data. Future empirical studies should therefore calibrate both the fractional order and the dissipation coefficient using operational observations collected from real logistics systems.

Second, the proposed formulation represents system dynamics through a single effective inventory-service state. While this aggregation preserves mathematical tractability, it neglects richer physical representations in which inventory, congestion, backlog, service quality, and product deterioration evolve as distinct but interacting state variables. Extending the framework toward higher-dimensional state representations would improve modeling realism for complex transportation networks.

Third, the carbon-trading mechanism assumes deterministic buying and selling prices. In practice, carbon markets are characterized by stochastic prices, evolving regulations, banking and borrowing mechanisms, and dynamic market conditions. Incorporating endogenous or stochastic carbon-price dynamics would transform the current operational model into an integrated operational-financial optimization framework and would better represent real cap-and-trade systems.

From an uncertainty perspective, the current formulation adopts a scenario-based stochastic programming approach. When historical information is limited or probability distributions are difficult to estimate reliably, robust or distributionally robust optimization may provide a more suitable modeling paradigm. Integrating ambiguity sets or robust uncertainty sets while preserving fractional state dynamics therefore constitutes an important direction for future research.

Another promising extension concerns demand forecasting. Rather than generating scenarios independently of historical observations, forecasting techniques such as ARIMA, SARIMA, Prophet, LSTM, or Transformer-based models could be integrated into a rolling-horizon optimization framework. Such an architecture would naturally interact with the fractional-memory formulation because forecasting errors would propagate through the historical state, allowing transportation decisions to account simultaneously for prediction uncertainty and long-memory effects.

The present model also assumes a common fractional order for all customers and transportation modes. In practical logistics systems, persistence may vary according to product characteristics, transportation technologies, customer profiles, or geographical regions. Future research may therefore consider customer-specific, mode-specific, or variable-order fractional operators to capture heterogeneous memory effects more realistically, although this would require more sophisticated calibration procedures.

From a sustainability perspective, another natural extension consists of reformulating the problem as a genuine multi-objective optimization model in which transportation cost, carbon emissions, service quality, and resilience are optimized simultaneously. Such a formulation would enable the construction of Pareto frontiers and provide deeper insights into the trade-offs introduced by fractional-memory dynamics.

The proposed framework could also be integrated into digital-twin architectures for intelligent transportation systems. Real-time operational data could continuously update both the system state and the uncertainty model, enabling adaptive transportation planning under dynamic operating conditions. Because fractional dynamics explicitly preserve operational history, they are particularly well suited to digital-twin environments that continuously synchronize physical and virtual systems.

From a computational viewpoint, further improvements are also desirable. The dense temporal coupling introduced by the L1 discretization increases computational complexity as the planning horizon and network size grow. Future work may investigate accelerated fractional convolution techniques, compressed-memory approximations, Benders decomposition, branch-and-cut algorithms exploiting the carbon-trading structure, or hybrid decomposition methods combining Progressive Hedging with Benders decomposition. More comprehensive computational studies including repeated SAA replications, confidence intervals, detailed convergence analyses, and scalability experiments with respect to the planning horizon, network size, and scenario set would further strengthen the practical applicability of the proposed methodology.

Overall, this work should be viewed as a methodological foundation for a broader research agenda at the intersection of fractional calculus, stochastic optimization, sustainable transportation, and environmental regulation. By explicitly incorporating long-memory dynamics into transportation planning, the proposed framework extends classical optimization models and opens new opportunities for integrating forecasting, robust optimization, digital twins, carbon-market modeling, and multi-objective sustainability within a unified decision-support framework.